% =============================================================================
% Template LaTeX IPG  —  https://github.com/VagnerBomJesus/TEMPLATE-IPG-LATEX
% Autor base, criador e mantenedor do projeto:
%     Vagner Bom Jesus  (@VagnerBomJesus)
%
% Proprietario e autor original deste template. Ao reutilizar, distribuir ou
% editar, MANTER este credito ao autor base (nao remover). Quem usa o template
% e o autor do SEU relatorio; o criador base do modelo e Vagner Bom Jesus.
% Licenca: GPL-3.0 (ver LICENSE).
% =============================================================================

% =============================================================================
% Capítulo 5 — Conclusões
% =============================================================================

\chapter{Conclusões}
\label{cap:conclusoes}

% Responder às seguintes questões:
% - Qual é o resumo do projeto (problema e solução)?
% - Qual é o resumo dos testes e resultados?
% - Quais são as limitações do trabalho?
% - O que ficou por fazer?
Este capítulo apresenta as conclusões do trabalho, as principais contribuições, as limitações identificadas e sugestões de trabalho futuro.

% ---- 5.1 Síntese do Trabalho ----
\section{Síntese do Trabalho}
\label{sec:sintese}

% Resumir o projeto: o problema abordado, a solução implementada e os
% principais resultados dos testes.
Apresentar uma síntese do trabalho realizado, recordando os objetivos propostos e os resultados alcançados.

% ---- 5.2 Contribuições ----
\section{Contribuições}
\label{sec:contribuicoes}

Identificar as principais contribuições do trabalho.

\begin{itemize}
    \item Contribuição 1
    \item Contribuição 2
    \item Contribuição 3
\end{itemize}

% ---- 5.3 Limitações ----
\section{Limitações}
\label{sec:limitacoes}

% Descrever o que o sistema não faz ou faz de forma limitada,
% e que pode ser melhorado no futuro.
Identificar as limitações do trabalho.

% ---- 5.4 Trabalho Futuro ----
\section{Trabalho Futuro}
\label{sec:trabalho_futuro}

% Descrever o que ficou por fazer, melhorias possíveis e
% direções de investigação futura.
Sugerir possíveis direções para trabalho futuro.

\begin{itemize}
    \item Melhoria 1
    \item Funcionalidade adicional 1
    \item Investigação futura 1
\end{itemize}
